\chapter{Ankle Sprain}

\section*{Definition and Overview}
An ankle sprain is an injury to one or more of the ligaments around the ankle joint, caused when the ankle is forced beyond its normal range of movement. Ligaments are strong bands of tissue that connect bone to bone and hold the joint stable. The large majority of sprains involve the lateral (outer) ligament complex, most often the anterior talofibular ligament, and occur when the foot rolls inwards (inversion).

Sprains are graded by severity: a grade 1 sprain is a mild stretching of the ligament with no significant tearing, a grade 2 sprain involves partial tearing with moderate pain and swelling, and a grade 3 sprain is a complete rupture of the ligament with marked swelling, bruising, and instability. Most sprains heal well with self-care and rehabilitation, but severe sprains can take months to recover and may leave the ankle feeling unstable.

\section*{Epidemiology}
Ankle sprains are among the most common musculoskeletal injuries in the world, accounting for a large proportion of emergency department attendances and of time lost from sport. They affect people of all ages but are most frequent in adolescents and young adults. Court sports such as basketball, netball, and tennis, together with soccer and other running and jumping sports, carry particular risk.

People who have sprained an ankle once are at substantially increased risk of doing so again, and inadequate rehabilitation is a major contributor to recurrent injury. Some studies show a slightly higher injury rate in women than in men, and older adults tend to sustain more severe injuries from comparable falls.

\section*{Aetiology and Pathophysiology}
The typical mechanism is the foot rolling inwards, which stretches the ligaments on the outside of the ankle beyond their elastic limit:

\begin{itemize}
    \item \textbf{Inversion injuries:} the most common type, injuring the lateral ligament complex
    \item \textbf{Eversion injuries:} the foot rolls outwards, straining the medial (deltoid) ligament; these are usually more serious
    \item \textbf{High-energy twists:} falls, uneven ground, and awkward landings from a jump
    \item \textbf{Sports injuries:} sudden changes of direction, landing on another player's foot, or stepping into a hole
    \item \textbf{Underlying factors:} previous sprains, weak muscles, poor proprioception and balance, and unsuitable footwear
\end{itemize}

When the ligament is stretched or torn, small blood vessels rupture, producing swelling and bruising, and the joint becomes painful and less stable. In a grade 3 sprain the ligament is completely torn, and the ankle may move abnormally when stressed.

\section*{Clinical Features}
Symptoms typically develop at the moment of injury:

\begin{itemize}
    \item A sudden pain, sometimes with a popping or tearing sensation
    \item Rapid swelling over the side of the ankle
    \item Bruising that may spread over the ankle and down into the foot over the following days
    \item Tenderness when pressing over the injured ligament
    \item Pain with weight-bearing or walking
    \item Stiffness and reduced range of movement
    \item A feeling that the ankle is unstable or wants to give way, more common in severe sprains
\end{itemize}

\section*{Differential Diagnosis}
Because many ankle injuries look alike, other conditions must be excluded:

\begin{itemize}
    \item \textbf{Ankle fracture:} a broken bone, diagnosed with X-ray
    \item \textbf{Syndesmosis injury:} a "high ankle sprain" involving the ligaments above the joint, which is usually slower to heal
    \item \textbf{Tendon injury:} tearing of the peroneal tendons or the Achilles tendon
    \item \textbf{Growth plate injury:} in children and adolescents, the growth plate may be injured instead of the ligaments
    \item \textbf{Bone bruise or osteochondral injury:} damage to the cartilage and bone at the joint surface
    \item \textbf{Tarsal or metatarsal fracture:} a break in the bones of the mid-foot or forefoot
\end{itemize}

\section*{When to Seek Medical Care}
Most mild sprains can be managed at home, but assessment is recommended whenever there is doubt. Seek care if:

\begin{itemize}
    \item You cannot put weight on the ankle for more than a few steps
    \item The pain is severe or the ankle looks deformed
    \item Swelling and bruising are significant or worsening
    \item Numbness, tingling, or a cold, pale foot develops
    \item You have had repeated sprains or the ankle keeps giving way
    \item Symptoms do not start to improve within a few days
\end{itemize}

\textbf{Call 000 (or local emergency services) or attend an emergency department for:}
\begin{itemize}
    \item A bone that is obviously out of place
    \item A cold, blue, or numb foot below the ankle
    \item A major mechanism of injury, such as a fall from a height or a car accident
\end{itemize}

\section*{Diagnosis and Investigations}
A doctor will ask how the injury happened and examine the ankle, checking for tenderness over the bones and ligaments:

\begin{itemize}
    \item \textbf{Ottawa ankle rules:} a widely used checklist that helps determine whether an X-ray is needed
    \item \textbf{X-ray:} the main test, used to exclude a fracture when there is bony tenderness or an inability to bear weight
    \item \textbf{Stress tests:} gentle movement of the ankle by the examiner to assess ligament looseness
    \item \textbf{Ultrasound:} can demonstrate stretched or torn ligaments and tendons
    \item \textbf{MRI:} reserved for severe or persistent injuries, or when cartilage damage is suspected
    \item \textbf{Re-assessment:} a repeat examination or imaging a few days later if symptoms are not settling
\end{itemize}

\section*{Management}
Most sprains are managed without surgery:

\begin{itemize}
    \item \textbf{RICE in the first few days:} Rest, Ice, Compression, and Elevation to limit swelling and pain
    \item \textbf{Pain relief:} paracetamol, or a short course of non-steroidal anti-inflammatory medicines (NSAIDs) where appropriate
    \item \textbf{Progressive weight-bearing:} gentle, gradual loading as pain allows, rather than total inactivity
    \item \textbf{Ankle support:} a compression bandage, brace, or strapping during early recovery
    \item \textbf{Physiotherapy:} exercises to restore movement, strength, and balance; this is the most important part of rehabilitation
    \item \textbf{Gradual return to activity:} a structured return to walking, running, and then sport, guided by symptoms
    \item \textbf{Surgery:} rarely needed, usually reserved for elite athletes or for severe sprains that fail to heal and leave the ankle persistently unstable
\end{itemize}

\section*{Complications}
\begin{itemize}
    \item \textbf{Chronic ankle instability:} ongoing looseness and repeated sprains after severe ligament tears
    \item \textbf{Persistent pain or stiffness:} occasionally lasting for months
    \item \textbf{Post-traumatic arthritis:} can develop years after repeated or severe injuries
    \item \textbf{Synovitis or scar tissue:} inflammation of the joint lining causing persistent swelling
    \item \textbf{Compartment syndrome:} a rare emergency in which severe swelling compresses nerves and compromises blood flow
    \item \textbf{Complex regional pain syndrome:} an uncommon reaction causing disproportionate pain, swelling, and skin changes
    \item \textbf{Secondary injuries:} an altered walking pattern can strain the foot, knee, hip, or back
\end{itemize}

\section*{Prognosis and Outcomes}
The outlook is generally good. Mild sprains usually settle within a few days to two weeks, moderate sprains may take a month or more, and severe sprains several months. Full recovery of strength, balance, and proprioception often takes longer than the resolution of pain, and returning to sport too early raises the risk of re-injury. With proper rehabilitation, most people regain full function. A minority develop long-term instability or pain, particularly when the initial injury was severe or was not fully rehabilitated.

\section*{Prevention and Self-Care}
\begin{itemize}
    \item Strengthen the muscles around the ankle and train balance with exercises such as single-leg standing
    \item Warm up and cool down properly before and after sport
    \item Wear supportive, well-fitted footwear appropriate for the activity and surface
    \item Use ankle braces or strapping during sport if you have had previous sprains
    \item Take care on uneven ground, and fix home hazards such as loose rugs and poor lighting
    \item After an injury, complete the full rehabilitation programme rather than stopping when the pain goes
    \item Maintain general fitness and muscle strength to reduce overall injury risk
\end{itemize}

\section*{Sources and Further Reading}
The following sources provide reliable further information on ankle sprains:

\begin{itemize}
    \item NHS. \textit{Information on sprains and strains, including ankle sprains.} https://www.nhs.uk/conditions/sprains-and-strains/
    \item Mayo Clinic. \textit{Overview of sprained ankles, including diagnosis and recovery.} https://www.mayoclinic.org/diseases-conditions/sprained-ankle/symptoms-causes/syc-20353225
    \item MSD Manual. \textit{Ankle sprains in the consumer health section.} https://www.msdmanuals.com/home/injuries-and-poisoning/sprains-and-strains/ankle-sprains
    \item American Academy of Orthopaedic Surgeons. \textit{Patient education on ankle sprains and rehabilitation.} https://orthoinfo.aaos.org/
    \item MedlinePlus. \textit{Ankle sprain: causes, treatment, and prevention.} https://medlineplus.gov/ankleinjuriesanddisorders.html
\end{itemize}
